\section{Ten Simple Rules for Better Figures}
Ich habe den Artikel "Ten Simple Rules for Better Figures"\cite{simple_rules} gewählt. Der wissenschaftliche Artikel ist öffentlich zugänglich und wurde von PLOS Computational Biology (2014) veröffentlicht.

\smallskip
\noindent
Der Artikel behandelt das Thema, wie man gute Diagramme bzw. im Allgemeinen Daten gut darstellt. Ich habe die Struktur des Artikels übernommen und fasse jeden Unterpunkt kurz zusammen und gehe darauf ein, wie es sich auf unser Projekt bezieht.

\begin{paracol}{2}
%
\smallskip
\begin{itemize}
    \item 1: Kenne die Nutzer
    
    Man muss beachten, für wen die Tabellen / Daten gedacht sind. Vor allem, dass es für all die verschiedenen Nutzer unterschiedliche Anforderungen und Voraussetzungen gibt.

    $\rightarrow$ Für uns ergibt sich diese Information aus den Personas, die wir entwickeln. Mit gut gewählten Personas können wir die App so entwickeln, dass die Komplexität und der Aufbau an den Großteil der Nutzer optimal angepasst sind.
    %
    \item 2: Identifiziere die Nachricht
    
    Die "Geschichte", die zu den Daten geführt hat, ist oft komplex. Bei der Darstellung muss man sich bewusst sein, dass man oft nur einen kleinen Teil davon visualisieren kann / sollte.

    $\rightarrow$ Besonders wichtig, da wir eine breite Bandbreite an unterschiedlichen Daten visualisieren sollen. Es kann und darf nicht für alle Datensätze die gleiche Visualisierung verwendet werden und die Hintergründe der Daten müssen beachtet werden.
    %
    \item 3: Beachte das Medium
    
    Die Visualisierung muss an das Medium angepasst werden. Es macht einen großen Unterschied, ob es gedruckter Text oder ein Programm ist. Die Größe der Tabelle ist auch wichtig, es macht einen Unterschied, ob der Nutzer am Handy oder am PC ist.

    $\rightarrow$ Wir erstellen eine App, daher wissen wir, welches Medium verwendet wird. Hier müssen wir aber trotzdem beachten, dass nicht jede Form der Visualisierung gut / sinnvoll ist.
    %
    \item 4: Unterschriften sind ein Muss
    
    Der Titel fasst es zusammen, es macht keinen Sinn, einfach irgendwelche Tabellen anzuzeigen. 

    $\rightarrow$ Wir an der Uni auch oft genug gesagt und gilt auch für unser Projekt.
    %
    \item 5: Achtung vor Standardwerten
    
    Es gilt zu beachten, dass Software mit potenziell problematischen Standardwerten kommt. 

    $\rightarrow$ Extrem wichtig bei uns, wir werden definitiv fertige Softwarepakete verwenden.
    %
\switchcolumn
    %
    \item 6: Farben gut nutzen
    
    Farben/Farbtabellen sollten so gewählt werden, dass alles ersichtlich ist. Jegliche Änderungen und Abweichungen sollten nur mit gutem Grund gemacht werden.

    $\rightarrow$ Wichtig, da unser App-Design auch zu den Visualisierungen passen sollte und Punkt 9 auch zu beachten ist.
    %
    \item 7: Niemanden in die Irre führen
    
    Daten sollten richtig skaliert und angezeigt werden. Wird das nicht beachtet, könnte es zu einer fehlerhaften Darstellung / Interpretation der Daten kommen.

    $\rightarrow$ Auch Teil unserer Personas: Daten sollten auch nicht verwendet werden können, um falsche / manipulative Schlüsse ziehen zu können.

    \item 8: Vermeide “Chartjunk”
    
    Nicht zu viele Farben und Informationen in eine Tabelle stecken. Die Visualisierung soll nicht überladen und unübersichtlich werden.

    $\rightarrow$ Da wir sehr viele Datensätze visualisieren, müssen wir das ganz besonders beachten.

    \item 9: Information vor Schönheit
    
    Ein einfacher Merksatz, dass die Visualisierung in erster Linie Daten richtig und gut darstellen muss und man erst danach auf künstlerische Aspekte achten sollte.

    $\rightarrow$ Eine wichtige Regel bei der Erstellung der App, um den Fokus nicht aus den Augen zu verlieren.

    \item 10: Richtiges Werkzeug verwenden
    
    Es soll aus all den Softwarepaketen und Programmiersprachen die richtige gewählt werden. Alle haben Vor- und Nachteile und je nach Visualisierung eignen sich manche mehr oder weniger.

    $\rightarrow$ In der App sind wir ein wenig limitiert, aber es ist trotzdem wichtig, welche Softwarepakete wir verwenden.
\end{itemize}
%
\end{paracol}